\documentclass[a4paper,12pt]{article}

\title{Blocco 3}
\author{Leonardo Centazzo}
\date{\today}


% Pacchetti essenziali
\usepackage{amsmath, amssymb, amsthm}
\usepackage{enumitem}
\usepackage{geometry}

% Impostazioni pagina minime
\geometry{a4paper, margin=2cm}

% Stili per teoremi, proposizioni, ecc.
\theoremstyle{plain}
\newtheorem{teorema}{Teorema}
\newtheorem{proposizione}{Proposizione}
\newtheorem{lemma}{Lemma}
\newtheorem{corollario}{Corollario}

\theoremstyle{definition}
\newtheorem{definizione}[teorema]{Definizione}
\newtheorem{esempio}[teorema]{Esempio}
\newtheorem{esercizio}{Esercizio}
\newtheorem{osservazione}[teorema]{Osservazione}

\newtheorem*{warning}{Warning}

\theoremstyle{remark}
\newtheorem{nota}[teorema]{Nota}


\begin{document}

\maketitle


\begin{warning}

	In the following we will adopt the following conventions: 
	\begin{itemize}
		\item In each exercise, $P$ will denote the poset used.
		\item When something coincide when seen in $V$ or in $V[G]$ we we'll not specify where we are looking at it.
			So for example we we'll always write $\omega$ instead of $\omega^{V}$, $\omega^{V[G]}$.
		\item We won't specify if something is considered in $V$ or in $V[G]$ if it is clear from the context.
	\end{itemize}

\end{warning}


For the next problem you can use that for inaccessible $\lambda$, $Col(\omega_1,<\lambda)$ has the $\lambda-cc$ and 
if $G$ is a $Col(\omega_1,<\lambda)$-generic filter over the ground model $V$, then $V[G] \models \lambda = \omega_2$.

\begin{esercizio}

	Suppose that $\lambda$ inaccessible. Show that $Col(\omega_1,<\lambda)$ forces that CH holds.
	
\end{esercizio}

\begin{proof}

	Since $\omega < \lambda$, we have by inaccessibility of $\lambda$ that $|2^{\omega}| < \lambda$.
	Let $G$ be a generic filter for $P$. First we claim that, since $P$ is $\omega_1$-closed, we have  $(2^{\omega})^V = (2^{\omega})^{V[G]}$.
	\begin{itemize}
		\item $(2^{\omega})^{V} \subseteq (2^{\omega})^{V[G]}$: because $V \subseteq V[G]$.
		\item $(2^{\omega})^{V} \supseteq (2^{\omega})^{V[G]}$: Let $V[G] \models A \subseteq \omega$.
			We want to show that $A \in V$. Let $\dot{A}$ be a $P$-name over $V$ for $A$.
			Now consider $\langle p_n : n < \omega \rangle$, defined inductively as follows:
			\begin{enumerate}

				\item $n = 0$:
					\begin{itemize}
						\item if $V[G] \models 0 \in A$ take $p_0 \in G$ such that $p_0 \Vdash \check{0} \in \dot{A}$
						\item if $V[G] \models 0 \not \in A$ take $p_0 \in G$ such that $p_0 \Vdash \check{0} \not \in \dot{A}$
					\end{itemize}

				\item $n \Rightarrow n+1$:
					\begin{itemize}
						\item if $V[G] \models n+1 \in A$ take $p_{n+1} \in G, p_{n+1} \leq p_n$ such that 
							$p_{n+1} \Vdash \check{n+1} \in \dot{A}$

						\item if $V[G] \models n+1 \not \in A$ take $p_{n+1} \in G, p_{n+1} \leq p_n$ such that 
							$p_{n+1} \Vdash \check{n+1} \not \in \dot{A}$


					\end{itemize}

			\end{enumerate}
	\end{itemize}
			Now let $\overline{p} \in G, \overline{p} \leq p_n$ for all $n<\omega$.
			Such an element exists because one can consider 
			$p' = \bigcup_{n} p_n, D = \{q : q \leq p'\} \cup \bigcup_{n} \{q: q \perp p_n\}$.
			We claim that $D$ is dense. In fact, let $r \in P$. If $r \perp p_n$ for some $n < \omega$ we already have $r \in D$.
			So we can assume $r \not \perp p_n$ for all $n < \omega$.
			Define $\langle r_n : n< \omega \rangle$ as follows:
			\begin{enumerate}
				\item $n=0$: take $r_0 \leq r,p_0$
				\item $n \Rightarrow n+1$: 
					\begin{itemize}
						\item if $r_n \perp p_{n+1}$ then we have $r_n \in D, r_n \leq r$ so one can stop. 
						\item othewise take $r_{n+1} \leq r_n, p_{n+1}$.
					\end{itemize}
			\end{enumerate}

			If the sequence stopped we proved the density, so assume the sequence did not stop. In this case take 
			$r' = \bigcup_{n} r_{n}$. Then we have $r' \leq p_n$ for all $n < \omega$, which implies $r' \leq p'$.
			So we have $r' \in D, r' \leq r$, which proves the density of $D$.

			Since $D \in M$ is dense and $G$ is generic we have $G \cap D \not = \emptyset$.
			Let $\overline{p} \in D \cap G$. Since $p_n \in G$ for all $n < \omega$ it must be $\overline{p} \leq p_n$ for all $n < \omega$.
			So $\overline{p}$ is the element we were looking for.

			Now consider $B = \{n : \overline{p} \Vdash \check{n} \in \dot{A}\}$. Note that $B \in M$ since one can use the equivalent 
			internal definiton of forcing. We claim $B = A$.

			\begin{itemize}
				\item $A \subseteq B$:
					$n \in A \Rightarrow \exists r \in G \big( r \Vdash \check{n} \in \dot{A} \big) \Rightarrow 
					\overline{p} \Vdash \check{n} \in \dot{A}$ because otherwise $\overline{p},r \in G$ would be incompatible.
				\item $A \supseteq B$:
					$n \in B \Rightarrow \overline{p} \Vdash \check{n} \in \dot{A} \Rightarrow V[G] \models n \in A$.
			\end{itemize}

			This means that $\mathcal{P}(\omega) ^V = \mathcal{P}(\omega) ^{V[G]}$.
			Now observe that $|2^{\omega}| ^V \geq |2^{\omega}|^{V[G]}$ and since $|2^{\omega}|^V \in \lambda$ we get 
			$|2^{\omega}|^{V[G]} \in \lambda$. Furthermore from $V[G] \models \lambda = \omega_2$ we get $V[G] \models |2^{\omega}| < \omega_2$.
			Since $V[G] \models ZFC$ and $|2^{\omega}| \geq \omega_1$ is a theorem of ZFC, we also get 
			$V[G] \models |2^{\omega}| \geq \omega_1$.
			So we conclude $V[G] \models |2^{\omega}| = \omega_1$.

\end{proof}



For the next problem, recall that $Add(\omega_1,1)$ is the poset to add one Cohen subset of $\omega_1$. More precisely,
$Add(\omega_1,1) = \{p: \omega_1 \to 2 : |p| \leq \omega\}$, ordered by reverse inclusion.

\begin{esercizio}

	Suppose that $G$ is $Add(\omega_1,1)$-generic over the ground model $V$. Show that $V[G] \models CH$.

\end{esercizio}

\begin{proof}
	
	As in the previous exercise, observe that $(2^{\omega})^V = (2^{\omega})^{V[G]}, (\omega_1)^V = (\omega_1)^{V[G]}$ because $P$ is $\omega_1$-closed.
	Now let $g = \bigcup G \in V[G]$. Then $g$ will be the characterictic function of a subset of $\omega_1$. Let's call it $A$.
	Now consider $f: \omega_1 \to \mathcal{P}(\omega), \alpha \mapsto B_{\alpha}$, where $n \in B_{\alpha} \iff \alpha + n \in A$.
	We claim $f$ surjective. In fact let $B \subseteq \omega$.We want to find $\alpha < \omega_1$ such that $B_{\alpha} = B$.
	Equivalently, we want to prove that $\{\alpha < \omega_1 : \forall n< \omega \big( n \in B \iff \alpha + n \in A \big) \} \not = \emptyset$.
	Consider $D_B = \{p \in P : \exists \alpha< \omega_1 \forall n<\omega 
	\bigg( \alpha + n \in dom(p) \land \big( p(\alpha + n) = 1 \iff n \in B \big) \bigg) \}$.
	Then $D_{B} \in M$ and $D_B$ is dense by regularity of $\omega_1$.
	Let $r_B \in D_B \cap G$. Then $r_B \Vdash \dot{B} \in \dot{ran(f)}$, which means that $V[G] \models f \text{ surjective }$.
	Equivalently, $V[G] \models CH$.

\end{proof}


\end{document}

